\chapter{Data Model and Persistence}

\section{Introduction}

All results produced by the pipeline are persisted so that they can be retrieved later through the API without reprocessing the video. The data layer is defined with SQLAlchemy~2.0 model definitions. Most tables carry a \code{video_id} that links them to the \code{Video} row they were derived from, which is why deleting a video (\code{DELETE /videos/{id}}) removes all associated data.

\section{Entities}

Table~\ref{tab:models} lists the persisted models and their key fields.

{\singlespacing\small
\begin{longtable}{@{}p{3.8cm}p{10.2cm}@{}}
\caption{Data models and key fields.}\label{tab:models}\\
\toprule
\textbf{Model} & \textbf{Key fields}\\
\midrule
\endfirsthead
\toprule
\textbf{Model} & \textbf{Key fields}\\
\midrule
\endhead
\bottomrule
\endfoot
\code{Video} & \code{id}, \code{filename}, \code{width}, \code{height}, \code{fps}, \code{frame_count}, \code{status}\\
\code{ProcessingJob} & \code{id}, \code{video_id}, \code{status}, \code{stage}, \code{progress}, \code{config} (dict)\\
\code{Experiment} & \code{id}, \code{name}, \code{tracker}, \code{scenario}, \code{num_frames}, \code{status}, \code{metrics}\\
\code{IdentityRecord} & \code{id}, \code{video_id}, \code{identity_id}, \code{prototype} (embedding vector), \code{quality}, \code{observation_count}, \code{is_confident}\\
\code{Track} & \code{id}, \code{video_id}, \code{track_id}, \code{frames}, \code{first_seen}, \code{last_seen}\\
\code{DetectionRecord} & \code{id}, \code{video_id}, \code{frame_id}, \code{bbox}, \code{confidence}\\
\code{FrameTimingRecord} & \code{id}, \code{video_id}, \code{frame_id}, \code{proc_ms}\\
\code{CameraMotionRecord} & \code{id}, \code{video_id}, \code{motion_type}, \code{compensation_applied}\\
\code{IdentityObservation} & \code{id}, \code{video_id}, \code{identity_id}, \code{track_id}, \code{frame_id}\\
\code{ReidEvent} & \code{id}, \code{video_id}, \code{lost_frame_id}, \code{reappearance_frame_id}\\
\code{AppSettings} & \code{id} (=\code{GLOBAL_SETTINGS_ID}), \code{values} (dict)\\
\end{longtable}}

\section{Entity Relationships}

Figure~\ref{fig:er} gives an overview of how the models relate to each other. \add{The relationships shown are inferred from the \code{video_id}, \code{identity_id}, \code{track_id} and \code{frame_id} fields listed above; \code{Experiment} and \code{AppSettings} are independent of any single video.}

\begin{figure}[htbp]
\centering
\begin{tikzpicture}[
  node distance=7mm and 8mm,
  ent/.style={draw,rounded corners=2pt,align=center,font=\footnotesize,minimum height=7mm,minimum width=2.6cm},
  main/.style={ent,fill=gray!15,font=\small\bfseries},
  lnk/.style={thick}]
\node[main] (video) {Video};
\node[ent,above left=of video] (job) {ProcessingJob};
\node[ent,above=of video] (track) {Track};
\node[ent,above right=of video] (det) {DetectionRecord};
\node[ent,left=of video] (ident) {IdentityRecord};
\node[ent,right=of video] (time) {FrameTimingRecord};
\node[ent,below left=of video] (obs) {IdentityObservation};
\node[ent,below=of video] (reid) {ReidEvent};
\node[ent,below right=of video] (cam) {CameraMotionRecord};
\node[ent,below=12mm of reid,xshift=-22mm] (exp) {Experiment};
\node[ent,below=12mm of reid,xshift=22mm] (set) {AppSettings};
\foreach \n in {job,track,det,ident,time,obs,reid,cam} \draw[lnk] (video) -- (\n);
\draw[lnk,dashed] (ident) -- (obs);
\draw[lnk,dashed] (obs) -- (track);
\end{tikzpicture}
\caption[Overview of the data model]{Overview of the data model: all per-video artefacts reference \code{Video}; \code{Experiment} and \code{AppSettings} are standalone.}
\label{fig:er}
\end{figure}

\section{Rows Created by One Processing Run}

When a processing job completes, the \code{Video} status becomes \code{processed} and the following rows are created: \code{IdentityRecord}, \code{Track}, \code{DetectionRecord}, \code{FrameTimingRecord} and \code{CameraMotionRecord}; identity sightings and re-identification events are stored as \code{IdentityObservation} and \code{ReidEvent} rows. Table~\ref{tab:rows} gives the template used to report the volume of data produced for a video.

\begin{table}[htbp]
\centering\small
\begin{tabular}{@{}lcc@{}}
\toprule
\hr \textbf{Model} & \textbf{Rows per video} & \textbf{Growth driver}\\
\midrule
\hr \code{DetectionRecord} & \pp{XX} & Frames $\times$ detections per frame\\
\hr \code{FrameTimingRecord} & \pp{XX} & One per processed frame\\
\hr \code{Track} & \pp{XX} & Number of tracks started\\
\hr \code{IdentityRecord} & \pp{XX} & Number of distinct fruit\\
\hr \code{IdentityObservation} & \pp{XX} & Sightings of identities\\
\hr \code{ReidEvent} & \pp{XX} & Occlusion/re-appearance events\\
\hr \code{CameraMotionRecord} & \pp{XX} & Motion estimate(s) for the video\\
\bottomrule
\end{tabular}
\caption[Rows per processed video]{\add{Rows created per processed video (counts to be filled in from a real run; the growth drivers are inferred).}}
\label{tab:rows}
\end{table}

\section{Identity Records}

An \code{IdentityRecord} stores the learned prototype of one fruit: an embedding vector, a quality score, the number of observations that contributed to it, and an \code{is_confident} flag. Each individual sighting of an identity in a frame, attributed to a track, is stored as an \code{IdentityObservation}. A \code{ReidEvent} records the frame in which an identity was lost and the frame in which it reappeared.

\add{This split keeps the compact summary (the prototype) separate from the evidence that produced it (the observations), which allows the interface to show both the identity-space visualisation and the per-frame history of any identity.}

\section{Tracks and Detections}

A \code{Track} holds the frames over which a tracker followed an object, together with its first and last frame. \code{DetectionRecord} stores the bounding box and confidence of each persisted detection. In the shared TypeScript types the bounding box is a \code{BoundingBox} with \code{x}, \code{y}, \code{width} and \code{height}, and a track point can be in one of three visibility states: \code{visible}, \code{occluded} or \code{lost}.

\section{Telemetry and Camera Motion}

\code{FrameTimingRecord} stores the processing time per frame in milliseconds (\code{proc_ms}) and is returned by \code{GET /videos/{id}/per-frame}. \code{CameraMotionRecord} stores the type of motion estimated for the video and whether compensation was applied, and is returned by \code{GET /videos/{id}/motion}.

\section{Settings and Experiments}

\code{AppSettings} is a single-row table whose \code{values} dictionary holds the pipeline defaults. \code{Experiment} stores the name, tracker, scenario, number of frames, status and computed metrics of an evaluation run. An experiment that is in the \code{RUNNING} state cannot be deleted.

\section{Database Selection}

SQLite is used throughout, with the development database file at \code{data/orchardvision.db} and the container stack pointing \code{DATABASE_URL} at a file in a shared volume. \add{Because all access goes through the SQLAlchemy ORM, the application code is identical whichever server is selected; only the connection string changes. The current limits of SQLite become visible with several concurrent writers, which is the main reason a deployment at scale would move to a server database. Every identifier is a 32-character hexadecimal UUID, and only the video-to-job relation is mapped as an ORM relationship; the remaining foreign keys are plain indexed columns, which keeps the persistence of high-volume tables such as detections and frame timings cheap.}
